\section{Classical Pfaff--Darboux theorems}
Let $\omega$ be a smooth $1$-form on an $n$-manifold~$M$
that, for some integer $k>0$, satisfies
\begin{equation}\label{eq: omegaPfaff_k}
\text{$\omega\w(\d\omega)^k$ vanishes identically on $M$}
\end{equation}
while
\begin{equation}\label{eq: omega_nondegenerate}
\text{$\omega\w(\d\omega)^{k-1}$ is nowhere vanishing on $M$.}
\end{equation}
The integer~$k{-}1$ is known as the \emph{Pfaff rank}
of~$\omega$~\cite[Chapter II, Section~3]{MR1083148}.
Note that $k\le \tfrac12(n{+}1)$. When $k=\tfrac12(n{+}1)$,
$\omega$ is said to be a \emph{contact form} on~$M$.

When $\omega$ satisfies~\eqref{eq: omegaPfaff_k}
and~\eqref{eq: omega_nondegenerate}, so does $\tilde\omega=f\omega$
for any nonvanishing function~$f$ on~$M$,
since $\tilde\omega\w(\d\tilde\omega)^{r-1} = f^r \omega\w(\d\omega)^{r-1}$
for all integers~$r>0$.

\subsection{Canonical subbundles}
An $\omega$ satisfying~\eqref{eq: omegaPfaff_k}
and~\eqref{eq: omega_nondegenerate} defines a \emph{kernel subbundle}~$K
= \omega^{-1}(0)\subset TM$ of corank~$1$
and a subbundle $A\subset K$ of corank $2(k{-}1)$ in~$K$
by the rule that, for each $m\in M$,
\begin{equation*}
A_m = \{ v\in K_m\mid \mathrm{d}\omega(v,w) = 0,\, \forall w\in K_m \}.
\end{equation*}
Replacing~$\omega$ by $\tilde\omega= f \omega$ for any nonvanishing
function~$f$ does not change~$K$ or $A$.

If $k>1$, then $K\subset TM$ is not an integrable plane field,
but the subbundle $A\subset TM$ is always integrable,
since it is the Cauchy characteristic plane field
of the differential ideal~$\cI$ generated by~$\omega$
(see~\cite[Chapter II, Proposition~2.1]{MR1083148}).
In the contact case, i.e., when $n=2k{-}1$ (which is,
in some sense, generic), one has~$A=(0)$.

There is a nondegenerate, skew-symmetric bilinear pairing
$B_\omega \colon K{/}A\times K{/}A\to\bbR$
defined by
\[
B_\omega\bigl(v{+}A_m,w{+}A_m\bigr) = \mathrm{d}\omega(v,w),
\]
when $v,w\in K_m$. It satisfies $B_{f\omega} = fB_\omega$
for any nonvanishing function~$f$ on~$M$.

Note that any subspace~$W\subset T_m$ on which both $\omega$ and $\d\omega$
vanish, must, first of all, satisfy $W\subset K_m$
(since $\omega$ vanishes on~$W$), and, second, must have codimension
at least $k{-}1$ in $K_m$, since $\d\omega$, as a skew-symmetric form
on $K_m$, has Pfaff rank~$k{-}1$.
Moreover, if $W$ does have codimension $k{-}1$ in $K_m$,
then it must contain $A_m$, so that $W/A_m$ is a null subspace of $B_\omega$.

\subsection{Legendrian submanifolds and Grassmannians}
Any submanifold~$L\subset M$ that satisfies~$L^*\omega=0$,
i.e., an \emph{integral manifold of~$\omega$},
must also satisfy $L^*\d\omega=0$
and hence, by the above linear algebra discussion,
must have codimension at least~$k$ in~$M$.
When $L\subset M$ is an integral manifold of~$\omega$ of codimension~$k$,
it is said to be an \emph{$\omega$-Legendrian submanifold}.

In particular, if $L\subset M$ is $\omega$-Legendrian,
then, for each $m\in L$, the tangent space $T_mL$
satisfies $A_m\subset T_mL\subset K_m$
while $B_\omega$ vanishes identically on $T_mL/A_m\subset K_m/A_m$.

This motivates defining
the \emph{Legendrian Grassmannian}~$\Leg_m(\omega)\subset\Gr^k(T_mM)$
to be the set of subspaces~$W\subset K_m$
that have codimension~$k$ in $T_mM$
and on which both~$\omega$ and~$\d\omega$ vanish.
By the above remarks, it follows that $\Leg_m(\omega)$
can be canonically identified with the
\emph{Lagrangian Grassmannian}~$\Lag(K_m/A_m)\subset\Gr^{k-1}(K_m/A_m)$
consisting of the $(k{-}1)$-dimensional subspaces of~$K_m/A_m$
on which $B_\omega$ vanishes.
Hence, $\Leg_m(\omega)$ is naturally a smooth manifold
of dimension~$\tfrac12k(k{-}1)$.
Moreover, the (disjoint) union
\[
\Leg(\omega) = \bigcup_{m\in M} \Leg_m(\omega)\subset \Gr^k(TM)
\]
is a smooth subbundle,
and $\Leg(f\omega)=\Leg(\omega)$ for all nonvanishing~$f$.

\subsection{A local normal form}
One version of the \emph{Pfaff--Darboux theorem}
\cite[Chapter II, Theorem~3.1]{MR1083148} states that,
when ${\omega\in\Omega^1(M)}$ satisfies~\eqref{eq: omegaPfaff_k}
and~\eqref{eq: omega_nondegenerate},
each $m\in M$
has an open neighborhood~$U\subset M$
on which there exist smooth functions~$u^1,\ldots,u^k$
and $a_1,\ldots,a_k$ (with not all $a_i$ simultaneously vanishing)
so that
\begin{equation}\label{eq: PfaffDarboux_normalform}
U^*\omega = a_1\,\d u^1 + \cdots + a_k\,\d u^k .
\end{equation}
Moreover,
the mapping~$(u,[a])\colon U\to\bbR^k\times\mathbb{RP}^{k-1}$
is a submersion. (In fact, the kernel subbundle of the differential
of this mapping is the restriction of~$A$ to~$U$.)

Conversely, the existence of functions $u^i$ and $a_i$ for~$1\le i\le k$
on an open set~$U\subset M$ satisfying~\eqref{eq: PfaffDarboux_normalform}
with the $a_i$ not all simultaneously vanishing and having the
property that $(u,[a]) \colon U\to\bbR^k\times\mathbb{RP}^{k-1}$
be a submersion implies that both~\eqref{eq: omegaPfaff_k}
and~\eqref{eq: omega_nondegenerate} hold on~$U$.

\subsection{Geometry of the normal form}
It will be useful to have a geometric interpretation
of the Pfaff--Darboux theorem.
Now, in the representation~\eqref{eq: PfaffDarboux_normalform},
the functions~$u^i$ have independent differentials,
i.e., $\d u^1\w\cdots\w\d u^k$ does not vanish on~$U$.
Consequently, the simultaneous level sets of the functions~$u^i$
define a foliation~$\cL$ of~$U\subset M$
by $\omega$-Legendrian submanifolds, i.e., an $\omega$-Legendrian foliation.

Conversely, given an $\omega$-Legendrian foliation~$\cL$
on an open subset~$V\subset M$, each point~$m\in V$ will
have an open neighborhood~$U\subset V$ in which the leaves of~$\cL$
are the fibers of a submersion~${u = (u^i) \colon U\to\bbR^k}$.
Since~$\omega$ vanishes when pulled back to any fiber of~$u$,
it follows that there exists a mapping~$a = (a_i) \colon U\to\bbR^k$
such that $U^*\omega = a_1\,\d u^1 + \cdots + a_k\,\d u^k$.

Thus, a geometric interpretation of the Pfaff--Darboux theorem
is the statement that,
when $\omega\in\Omega^1(M)$ satisfies~\eqref{eq: omegaPfaff_k}
and~\eqref{eq: omega_nondegenerate}, each point~$m\in M$
has an open neighborhood~$U\subset M$
on which there exists an $\omega$-Legendrian foliation.

